% !PRINTED 607
\weberpage{607}

\begin{center}
\textbf{\S\,164.\qquad \gesperrt{Der Klassenkörper.}}
\end{center}

Zu der Zahlgruppe $A$ in dem quadratischen Körper $\Omega$ soll ein algebraischer Körper über $\Omega$ existieren, den man mit $\mathfrak{K}$ oder $\mathfrak{K}(A)$ bezeichne, von dem wir nur voraussetzen wollen:

7. \gesperrt{Definition des Klassenkörpers.}

\gesperrt{Die Primideale $\mathfrak{p}_1$ ersten Grades der Hauptklasse $\overline{A}_1$, und nur diese, sollen im Körper $\mathfrak{K}(A)$ wieder in Primideale ersten Grades zerfallen.}

Es sind also nur die Primideale in $\Omega$, die wir vorhin mit $\mathfrak{p}_1$ bezeichnet haben, die in $\mathfrak{K}$ in Primideale ersten Grades zerfallen.

Eine beliebige endliche Anzahl von Primzahlen können dabei ausgenommen sein. Sie werden dann in den Exkludenten genommen.

Ist $n$ der relative Grad des Körpers $\mathfrak{K}$ über $\Omega$ und $p$ eine natürliche Primzahl, in der die Ideale $\mathfrak{p}_1$ und $\mathfrak{P}$ in $\Omega$ und $\mathfrak{K}$ aufgehen, und bezeichnet man mit $N_\Omega, N_\mathfrak{K}$ die absoluten Normen im Körper $\Omega$ und $\mathfrak{K}$, so ist
\[
N_\mathfrak{K}(\mathfrak{P}) = N_\Omega(\mathfrak{p}_1) = p, \quad N_\mathfrak{K}(\mathfrak{p}_1) = p^n,
\]
und folglich muß $\mathfrak{p}_1$ im Körper $\mathfrak{K}$ in $n$ Primideale $\mathfrak{P}$ zerfallen:
\weq{1}{\mathfrak{p}_1 = \mathfrak{P}_1 \mathfrak{P}_2 \ldots \mathfrak{P}_n.}
Die Primideale, die in den Grundzahlen der Körper $\Omega$ oder $\mathfrak{K}$ aufgehen, sollen immer ausgeschlossen sein. Dann sind die $\mathfrak{P}_1, \mathfrak{P}_2, \ldots, \mathfrak{P}_n$ voneinander verschieden.

Ob ein solcher Körper $\mathfrak{K}$ existiert, bleibt vorläufig unentschieden. Wir wollen untersuchen, was aus der Definition geschlossen werden kann.

Da in jedem Ideal $\mathfrak{p}_1$ $n$ Ideale $\mathfrak{P}$ aufgehen, so ist
\weq{2}{\prod\nolimits \left[1 - N_\mathfrak{K}(\mathfrak{p}_1)^{-s}\right]^n = \prod\nolimits \left[1 - N_\mathfrak{K}(\mathfrak{P})^{-s}\right],}
worin sich das erste Produkt auf alle Ideale $\mathfrak{p}_1$ des Körpers $\Omega$, das zweite auf alle Ideale $\mathfrak{P}$ ersten Grades des Körpers $\mathfrak{K}$ erstreckt, und nach dem für jeden beliebigen Körper gültigen Satz I, § 197 des II. Bandes hat
\weq{3}{(s-1)\prod\nolimits \frac{1}{\left[1-N_\mathfrak{K}(\mathfrak{P})^{-s}\right]}}
für $s=1$ einen von Null verschiedenen endlichen Grenzwert, der dort auf die Klassenzahl im Körper $\mathfrak{K}$ zurückgeführt ist. Setzen wir zur Abkürzung
%CONT>>
